\chapter{SOCI Ledger --- Sociology Tags}
\label{chap:SOCI_Ledger}

\section{Purpose}

The SOCI ledger classifies Engrams according to sociological structures,
group dynamics, collective meaning-making, and cultural context.

\section{Scope}

Tags in this ledger describe:
\begin{itemize}
    \item social roles and relationships,
    \item group identities,
    \item ritual and shared practices,
    \item institutions and cultural norms,
    \item societal pressures or expectations.
\end{itemize}

\section{Schema}

A SOCI tag is defined as:
\[
T_{\text{SOCI}} = \langle key,\, gloss,\, bindings,\, constraints \rangle
\]

Example tags:
\begin{itemize}
    \item \texttt{soci:role}
    \item \texttt{soci:ritual}
    \item \texttt{soci:collective-action}
    \item \texttt{soci:identity-group}
\end{itemize}

\section{Class Binding Notes}

\begin{itemize}
    \item O-class common, interacts with narrative and world-building.
    \item S-class allowed for structured societal systems.
    \item P-class only when representing moral/social commitments.
\end{itemize}

\section{Constraints}

\begin{itemize}
    \item Tags must not stereotype populations.
    \item Claims about groups must remain descriptive, not prescriptive.
    \item HITL needed for P-class social identity assertions.
\end{itemize}
